% !TEX root = ../main.tex
\clearpage
\chapter{Appendix}
\label{ch:appendix}

% Demote headings one level under Chapter 8 (sections become 8.1.x).
\renewcommand{\dissertationSubheading}[2][]{%
  \phantomsection
  \subsection{#2}%
  \if\relax\detokenize{#1}\relax\else\label{#1}\fi
}
\renewcommand{\dissertationSubsubheading}[2][]{%
  \phantomsection
  \subsubsection{#2}%
  \if\relax\detokenize{#1}\relax\else\label{#1}\fi
}
\section{Evaluation Pipeline and Reproducibility}
\label{ch:appendix-eval}

This appendix documents the evaluation pipeline behind Chapter~\ref{ch:results} and behind the registered analysis of Appendix~\ref{sec:map-registered}.

\dissertationSubheading[sec:pipeline-stages]{Pipeline stages}

The evaluation source is in the folder \texttt{experiments} of the public repository \url{https://github.com/abitocodes/ERC20-Allowance-PageRank-Wallet-Reputation}, and the paths below are relative to it. The evaluation has seven stages. The first four are those of the registered analysis: three run on Arbitrum One over the window of scan~A, 1~October~2023 to 30~June~2026 (UTC, inclusive at both ends), and the fourth adds July to September~2026. The last three build and evaluate the contract graph of the main analysis.

\begin{enumerate}
\item Plan, cohort and graphs: Commit the analysis plan (\texttt{docs/analysis\_plan.md}) and the configuration before any log of the study is extracted. Copy the logs of the window into a project table (scan~A), list the spenders that at least three owners approved (\texttt{sql/02\_spender\_candidates.sql}), read their account types (\texttt{scripts/check\_account\_code.py}) and keep the contracts as the plan's cohort (\texttt{scripts/build\_cohort.py}). Extract the ego events of the cohort (\texttt{sql/03\_ego\_events.sql}) and aggregate them into the pair tables of both graphs at $T_{\mathrm{obs}}$ and at $t_1$ (\texttt{sql/04\_anchor\_pairs.sql}).
\item Checks and alignment: Score every method at $T_{\mathrm{obs}}$ with the solver settings of Section~\ref{sec:notation}, compute the checks of Section~\ref{sec:proxy-formulas} and report Spearman $\rho$ and Kendall $\tau$ with 95\% paired-bootstrap intervals (400 resamples, seed 42) and the contrasts of Section~\ref{sec:statistical-tests} (\texttt{scripts/run\_same\_window.py}).
\item Temporal holdout W0: Score every method at $t_1$ and evaluate the spender cohort against the labels of April--June~2026, with the contrasts of rule~A (\texttt{scripts/run\_holdout.py}).
\item Registered window W1: Commit the registration (\texttt{docs/registration\_w1.md}) with the status in the configuration set to registered. \texttt{scripts/extract\_w1.py}, which will not run before that commit, runs scan~B and the ego extraction of July--September and writes an extraction manifest; \texttt{scripts/run\_registered.py} evaluates rule~B.
\item Tokens, prices, cohort lists and labels: \texttt{scripts/usd\_inputs.py} builds the token list and the daily prices, and \texttt{scripts/chainlink\_check.py} compares them with the Chainlink feeds. \texttt{scripts/run\_usd\_pipeline.py} and \texttt{scripts/run\_usd12m\_pipeline.py} load both and, from the five-token events of the ego tables of stages 1 and~4, list for each freeze the contract spenders that at least three owners approved in its score window and count the labels of the spenders with a positive latest allowance at the freeze (\texttt{sql/21}, \texttt{sql/16}).
\item Contract graph: \texttt{scripts/run\_contract12m\_pipeline.py} runs the BigQuery jobs (\texttt{sql/24} to \texttt{sql/43}) and the RPC reads of the contract graph in a fixed order, in two groups of steps that it calls gate~A and gate~B. The steps of gate~A are \texttt{check}, \texttt{subset}, \texttt{ends}, \texttt{activity}, \texttt{queue}, \texttt{typing}, \texttt{load}, \texttt{contracts}, \texttt{balances}, \texttt{anchors}, \texttt{audit} and \texttt{structure}: they read the five-token events of 1~April~2025 to 30~September~2026, type the accounts (\texttt{scripts/check\_account\_code.py}), keep the events between contracts, reconstruct and check the balances (\texttt{scripts/balance\_audit.py}) and end with the structure pilot of Section~\ref{sec:discovery-protocol}. The steps of gate~B are \texttt{constants}, \texttt{pairs}, \texttt{rest}, \texttt{checks} and \texttt{download}: they compute the capped allowances, the constants, the pair tables, the cohorts, the proxies, the label parts and the importance, check 64 invariants, each a count of violations that must be zero, and download the tables. The run manifest records every job, the result of every step and every decision by which the author let a step go on; the digest quotes these decisions.
\item Evaluation of the contract graph: With \texttt{CONTRACT\_REP\_VARIANT=contract12m} the evaluation scripts read the downloaded tables, in this order (all in \texttt{scripts/}): \texttt{run\_same\_window.py} (with the damping sweep and the secondary part of rule~A), \texttt{run\_discovery.py} (the four runs of the discovery-order test), \texttt{run\_holdout.py} (W0 and the primary part of rule~A), \texttt{run\_registered.py} (W1: F1, F2 and F6 with the two registered sensitivity analyses), \texttt{run\_robustness.py} (value slope, sample definition and the uncapped allowance layer) and \texttt{run\_benchmark.py}. Then \texttt{run\_supplementary.py} computes, from the saved scores and with the same resamples, the contrasts that belong to no rule (Table~\ref{tab:cpr-contrasts}), \texttt{sybil\_model.py} the quantities of Section~\ref{sec:sybil-cost-model} and \texttt{describe\_data.py} those of Appendix~\ref{ch:appendix-data}; \texttt{publish\_results.py} writes the published summaries with their digest, and \texttt{export\_latex.py} the tables in \texttt{results/tables/} and the figures in \texttt{Figures/generated/}.
\end{enumerate}

The timed runs were made with no other job of the pipeline running, on a machine with an Intel Core Ultra~9 185H processor and 64~GB of memory under Windows~11 (build 26200), with Python 3.11.9, NumPy 2.3.2 and SciPy 1.16.1.

The experiments folder holds the queries, scripts and tests, the versioned configuration (\texttt{config/contract\_reputation.yaml} with the overlay \texttt{config/overlay\_contract12m.yaml}), the plan, the registration and the four documents that define the main analysis (\texttt{docs/}: \texttt{revision\_price\_weighting.md}, \texttt{revision\_one\_year\_window.md}, \texttt{revision\_contract\_graph.md} and \texttt{discovery\_protocol.md}), the token list and daily prices with their Chainlink comparison (\texttt{data/0-usd-token-list-and-prices/}) and the cost ledger. The summaries behind the tables are those of commit \texttt{6b389df} of the public repository, in \texttt{data/3-published-results-for-thesis/revised-contract-12m/}, and those of the registered analysis are in \texttt{data/3-published-results-for-thesis/}; the digest of the main analysis records the code revision at which it was written. Each of the two folders holds copies of the configuration, the plan and the registration, a digest \texttt{results\_digest.md}, which names the summary file and key of every number, and \texttt{SHA256SUMS.txt}, so any number in a table can be checked against its summary without re-running the pipeline. The extracted tables and the per-contract scores, stored as parquet parts below 90\,MB, are kept in a public Google Drive folder; the committed file \texttt{drive-manifest.json} gives its link and lists every part with its checksum.

Synthetic tests check the solver, the edge weights, the paired bootstrap, the US-dollar conversion, the balance cap, the account typing and the crawl offline. Real extraction needs Google Cloud credentials and a billing project.

\dissertationSubheading[sec:bq-extraction]{BigQuery extraction design}

The public table \texttt{bigquery-public-data.goog\_blockchain\_arbitrum\_one\_us.logs} is read once for each window. Scan~A copies the ERC-20 \texttt{Approval} and \texttt{Transfer} logs from 1~October~2023 to 30~June~2026 into one project table partitioned by month, and scan~B does the same for July to September~2026. Every later query on these logs reads the project tables, so the chain-wide scan is paid for once per window. The registered analysis extracts from them the events with a contract of the plan's cohort at one end. The main analysis draws from them the events of the five tokens from 1~April~2025 to 30~September~2026 into one table, which its later jobs read: an edge needs a contract at both ends, so these events come from the whole chain and not from the neighborhood of a cohort. From the public data set it also reads the transaction senders among the addresses of these events, with two cross-checks of the typing, the last block at or before each freeze and the \texttt{Transfer} logs of the five tokens before 1~October~2023, with which the balances begin. Appendix~\ref{ch:appendix-sql} shows the queries.

Every job is listed in \texttt{data/bq\_ledger.json} with its query, parameters, estimated and billed bytes and cost. The jobs of the study billed $8.78$\,TiB, or USD~54.88 at the on-demand rate of USD~6.25 per TiB, against the budget of USD~60 set in the configuration of the main analysis. The jobs of the registered analysis billed USD~29.71 and the 26 jobs of the contract graph USD~12.63; the ledger lists the other jobs one by one, those of stage~5 among them.

\dissertationSubheading[sec:account-type-classification]{Account-type classification}

Section~\ref{sub:account-types} states the rule by which an address is a user account or a contract; this section documents the reads behind it. An address that sent a transaction is a user account without a read. \texttt{scripts/check\_account\_code.py} read every other address that can end an edge between contracts, every member of the cohort and spender lists and two audit samples of $1{,}000$ addresses each with \texttt{eth\_getCode} at block $513{,}216{,}506$, from the public endpoints \texttt{arb1.arbitrum.io} and \texttt{arbitrum-one.public.blastapi.io}. At that block the registered analysis had read the $29{,}022$ candidates of the plan's cohort, of which $27{,}844$ held contract code, $1{,}121$ no code and $57$ an EIP-7702 delegation. The block lies after the end of the data, so every address is typed by the code it held on that day.

No code, or a delegation designator (\texttt{0xef0100} followed by 20 bytes), means an externally owned account. For other code the script recognizes ERC-1167 clones \parencite{murray2018erc1167}, ERC-1967 proxies \parencite{palladino2019erc1967} and Safe proxies \parencite{safenddocs}, reads their implementation and calls \texttt{getThreshold()}, \texttt{getOwners()} and \texttt{entryPoint()}: a Safe answers a threshold from 1 to $1{,}000$ and a non-empty list of owners, and an ERC-4337 account \parencite{buterin2021erc4337} answers one of the four canonical EntryPoint contracts (versions 0.6 to 0.9) and holds, in its code or its implementation's, a \texttt{validateUserOp} selector and no \texttt{validatePaymasterUserOp} selector. The calls are made per address for ERC-1967 and Safe proxies and for code that answers \texttt{getThreshold()}, and once per code hash otherwise. Table~\ref{tab:contract-types} in Appendix~\ref{sec:data-contract-graph} gives the counts and the code families that complete the rule. No call was aborted at every attempt, and none of the $1{,}000$ audited transaction senders held contract code.

\dissertationSubheading[sec:benchmark-run-logs]{Benchmark run logs}

Table~\ref{tab:benchmark-runs} lists the wall-clock seconds of every timed solver call of the benchmark, in the order the runs were made, with their mean and standard deviation. Each configuration had one untimed warm-up before its five timed runs, and each graph and setting was timed once: the $d=0.85$ rows of Table~\ref{tab:robustness-damping} and the whole-graph row of Table~\ref{tab:benchmark-scaling} repeat the first block of this table.

% Generated by experiments/scripts/export_latex.py -- do not edit by hand.
% Source: experiments/data/2-processed-tables-and-evaluations/revised-contract-12m/benchmark/benchmark.json
\begin{table}[htbp]
\centering
\caption[Every timed solver call of the benchmark.]{Wall-clock seconds of every timed solver call of the benchmark (matched cohort, $n=6{,}443$ contracts), in the order the runs were made, with their mean and standard deviation. Each configuration had one untimed warm-up before its timed runs, and each graph and setting was timed once: the $d=0.85$ rows of Table~\ref{tab:robustness-damping} and the full-cohort row of Table~\ref{tab:benchmark-scaling} are the first block, reused. The stages are the subsamples of Table~\ref{tab:benchmark-scaling}.}
\label{tab:benchmark-runs}
\fitwidth{%
\begin{tabular}{lrrrrrrr}
\toprule
Method & Run 1 & Run 2 & Run 3 & Run 4 & Run 5 & Mean & SD \\
\midrule
\multicolumn{8}{l}{\emph{Full graphs, $d=0.85$}} \\
EndorseRank & 0.101 & 0.112 & 0.096 & 0.107 & 0.089 & 0.101 & 0.009 \\
AWP & 2.964 & 2.792 & 2.852 & 2.514 & 2.984 & 2.821 & 0.189 \\
C-PR ($\lambda=0.5$) & 2.954 & 2.963 & 3.394 & 3.095 & 2.993 & 3.080 & 0.184 \\
S-PR & 2.997 & 2.682 & 2.512 & 2.527 & 2.565 & 2.657 & 0.202 \\
\midrule
\multicolumn{8}{l}{\emph{Full graphs, $d=0.75$}} \\
EndorseRank & 0.082 & 0.088 & 0.093 & 0.085 & 0.081 & 0.086 & 0.005 \\
AWP & 2.146 & 1.687 & 1.613 & 1.562 & 1.769 & 1.755 & 0.232 \\
\midrule
\multicolumn{8}{l}{\emph{Full graphs, $d=0.95$}} \\
EndorseRank & 0.179 & 0.233 & 0.241 & 0.204 & 0.185 & 0.208 & 0.028 \\
AWP & 6.870 & 6.961 & 7.566 & 7.555 & 6.957 & 7.182 & 0.348 \\
\midrule
\multicolumn{8}{l}{\emph{Stage $n=1{,}000$}} \\
EndorseRank & 0.001 & 0.001 & 0.001 & 0.001 & 0.001 & 0.001 & 0.000 \\
AWP & 0.096 & 0.096 & 0.095 & 0.095 & 0.095 & 0.095 & 0.000 \\
\midrule
\multicolumn{8}{l}{\emph{Stage $n=2{,}000$}} \\
EndorseRank & 0.002 & 0.002 & 0.002 & 0.002 & 0.002 & 0.002 & 0.000 \\
AWP & 0.135 & 0.131 & 0.132 & 0.138 & 0.132 & 0.134 & 0.003 \\
\midrule
\multicolumn{8}{l}{\emph{Stage $n=5{,}000$}} \\
EndorseRank & 0.007 & 0.006 & 0.006 & 0.006 & 0.006 & 0.006 & 0.000 \\
AWP & 0.693 & 0.735 & 0.698 & 0.657 & 0.663 & 0.689 & 0.031 \\
\bottomrule
\end{tabular}
}
\end{table}


\dissertationSubheading[sec:parameter-summary]{Parameter summary}

Table~\ref{tab:parameter-summary} collects the default parameters of the main analysis; every value is read from the versioned configuration or, for the tokens, prices and constants, from the inputs committed with the analysis (Section~\ref{sec:amount-revision}), and $\rho$, $q$ and the primary reveal mode are those that the rules of the protocol set (Section~\ref{sec:discovery-protocol}).

\begin{table}[!htb]
\centering
\caption{Default parameters used in the evaluation.}
\label{tab:parameter-summary}
\fitwidth{%
\begin{tabular}{lll}
\toprule
Parameter & Value & Role \\
\midrule
PageRank damping $d$ & $0.85$; sweep $0.75$, $0.95$ & Teleport mixture \\
Convergence tolerance $\varepsilon$ & $10^{-8}$ in $L_1$ norm & Power-iteration stop \\
Max iterations $I_{\max}$ & $300$ & Power-iteration cap \\
Logistic $k$ & $0.01$ & Decay slope (AWP, EndorseRank, transfer layer of C-PR) \\
Logistic $t_0$ & $180$ days & Decay midpoint (AWP, EndorseRank, transfer layer of C-PR) \\
Value transform $b$ & $\ln 3/m$ per dollar, $m_T=160$, $m_A=0.071$ & Value weight $V(z)$ of AWP ($m_T$) and EndorseRank ($m_A$) \\
Balance cap & allowance counted up to the owner's balance at the freeze & Every allowance weight \\
Allowance cap $U$ & $16{,}000$ dollars & C-PR allowance layer; in-approve value \\
Cap per user account $U_{\mathrm{user}}$ & $31{,}000$ dollars & Importance $I_T$ and $I_A$ of the discovery-order test \\
Tokens and prices & WETH, USDC, USDT0, ARB, WBTC; DefiLlama daily close & Events and US-dollar amounts (Table~\ref{tab:usd-tokens}) \\
Restart distribution & $\propto X_u$; uniform over the contracts of a graph & AWP; EndorseRank, C-PR \\
Coupling weight $\lambda$ & $0.5$ primary; $0.25$, $0.75$ & C-PR allowance-layer weight, fixed in the plan \\
S-PR restart vector & C-PR ($\lambda=1$) on the transfer-layer nodes, renormalized & Endorsement-seeded teleportation \\
Score windows & 2025-07-01 -- 2026-06-30; 2025-04-01 -- 2026-03-31 & Same window and W1; W0. UTC inclusive; the freeze anchors the decay \\
Events of the contract graph & 2025-04-01 -- 2026-09-30 & Both score windows and both label windows \\
Account-type block & $513{,}216{,}506$ & Code and calls of the typing \\
Extraction (scan~A) & 2023-10-01 -- 2026-06-30 & Scored in full by the registered analysis \\
W0 freeze / labels & 2026-03-31 / 2026-04-01 -- 2026-06-30 & Spender cohort; rule~A \\
W1 freeze / labels & 2026-06-30 / 2026-07-01 -- 2026-09-30 & Spender cohort; F1 and F2 of rule~B, margin $\delta=0.01$ \\
Bootstrap resamples / seed & $400$ / $42$ & Paired bootstrap over contracts; over hot contracts in the crawls \\
Crawl seeds & $20$ contracts of the giant SCC & Starts of the crawls \\
Recompute share $\rho$ & $1/50$; $1/200$ on five seeds & Schedule of the orderings; sensitivity \\
Hot share $q$ / area share $x$ & $1\%$; $5\%$ / $10\%$; $5\%$, $20\%$ & Hot set $H_q(I)$ / $\mathrm{AUC}_x$; sensitivities \\
Reveal mode & \texttt{both}; \texttt{out} & Primary; sensitivity \\
Timed runs / warm-up & $5$ / $1$ & Runtime protocol \\
Scaling seed & \texttt{contract-benchmark-v1} & Deterministic SHA256 subsamples \\
BigQuery budget & USD~60 & Checked before every query \\
\bottomrule
\end{tabular}
}
\end{table}

% !TEX root = ../main.tex
\section{SQL Templates}
\label{ch:appendix-sql}

This appendix shows the two query patterns on which the contract graph rests, abridged; the full files, with those of the registered analysis, are in \texttt{sql/} of the experiments folder (Appendix~\ref{sec:pipeline-stages}). Names that start with \texttt{@} are query parameters.

\dissertationSubheading[sec:sql-materialize]{Event selection pattern}

Scans~A and~B tag each log of the public table by the topic hash of its event and store its emitter, its two indexed addresses \texttt{a} and \texttt{b} and its 32-byte amount (\texttt{sql/01\_materialize\_logs.sql}). \texttt{sql/24\_c12\_logs5.sql} keeps the events of the five tokens from both scan tables, and \texttt{sql/31\_c12\_contracts.sql} those between two contracts. Their filters, abridged:

\begin{verbatim}[baselinestretch=1]
-- sql/24: the events of the five tokens (src: the rows of both scan tables)
FROM src AS s
WHERE s.emitter IN (SELECT token FROM `dissertation-bq.contract_rep.u_tokens`)
  AND (s.kind = 'A' OR LTRIM(TO_HEX(s.val), '0') != '')
-- sql/31: the events between contracts
FROM `dissertation-bq.contract_rep.c12_logs5` AS l
WHERE l.a IN (SELECT address FROM `dissertation-bq.contract_rep.c12_contracts`)
  AND l.b IN (SELECT address FROM `dissertation-bq.contract_rep.c12_contracts`);
\end{verbatim}

For an approval (\texttt{kind} A), \texttt{a} is the owner and \texttt{b} the spender; for a transfer (T), \texttt{a} is the sender and \texttt{b} the recipient. Approvals of zero are kept, because an approval of zero sets an allowance to zero; a transfer is kept only when its value is not zero, and it is valued at its token's price on the UTC day of its block. \texttt{c12\_contracts} holds the addresses typed as contracts (Appendix~\ref{sec:account-type-classification}).

\dissertationSubheading[sec:sql-pairs]{Latest-allowance and pair-table pattern}

\texttt{sql/36\_c12\_allow\_triples.sql} reduces the approvals between contracts to the latest allowance of each token, owner and spender in the score window of a freeze and caps it at the owner's balance; \texttt{sql/37\_c12\_pairs.sql} sums the weights of each pair. Abridged, for $T_{\mathrm{obs}}$:

\begin{verbatim}[baselinestretch=1]
latest_tobs AS (                                  -- sql/36; latest_t1 likewise
  SELECT token, owner, spender, value, block_timestamp FROM (
    SELECT token, a AS owner, b AS spender, value, block_timestamp,
           ROW_NUMBER() OVER (PARTITION BY token, a, b
                              ORDER BY block_number DESC, log_index DESC) AS rn
    FROM `dissertation-bq.contract_rep.c12_events_cc`
    WHERE kind = 'A' AND block_timestamp >= TIMESTAMP(@tobs_start)
      AND block_timestamp <= TIMESTAMP(@tobs))
  WHERE rn = 1 AND value > 0),
-- with the decimals, the anchor day's price_usd and the owner's balance bal:
LEAST(value / POW(10.0, decimals) * price_usd,
      IF(bal >= POW(10.0, decimals - CAST(@dust_digits AS INT64)), bal, 0.0)
        / POW(10.0, decimals) * price_usd) AS usd_eff
-- sql/37, per (tag, owner, spender) over the triples with usd_eff > 0:
SUM(sig(age_days) * vt(usd_eff, CAST(@b_allowance AS FLOAT64))) AS w_er,
SUM(LEAST(usd_eff, CAST(@cap_allowance AS FLOAT64))) AS w_raw,
-- and per (tag, sender, recipient) over the transfers of c12_events_cc:
-- w_awp = SUM(sig * vt(usd, @b_transfer)), w_raw = SUM(usd * sig)
\end{verbatim}

The window function picks the latest approval of each token, owner and spender by block number and log index; it is kept only when its value is positive, so a revocation removes that token from the pair. \texttt{usd\_eff} is the smaller of the allowance and the owner's balance, both at the price of the anchor's day, with a balance below $10^{-6}$ token counted as zero (Section~\ref{sub:balances}). \texttt{w\_er} and \texttt{w\_awp} are the weights of EndorseRank and AWP and \texttt{w\_raw} those of the layers of C-PR; \texttt{sig} and \texttt{vt} are the logistic decay of Equation~\eqref{eq:logistic-decay} and the value transform of Equation~\eqref{eq:value-transform}, set to one from $bz=40$ on, where it equals one in double precision. Each pair row carries the weights of every method, so all scores at one anchor come from the same rows. The same query writes the weights with the slope multiplied and divided by ten, for the check of Table~\ref{tab:robustness-slope}, and the pairs without the balance cap, for that of Table~\ref{tab:uncapped-sensitivity}.

% !TEX root = ../main.tex
\section{Elementary Properties of PageRank and Rank Correlation}
\label{ch:appendix-theory}

Chapters~\ref{ch:literature}--\ref{ch:results} rely on the elementary properties summarized here. The material is standard: the PageRank properties follow \textcite{page1999pagerank} and \textcite{langville2006pagerank}, and the rank correlations \textcite{spearman1904proof} and \textcite{kendall1938rank}.

\dissertationSubheading[sec:pr-existence]{Existence and uniqueness of PageRank}

Take $P$ to be a row-stochastic matrix on $|V|$ nodes, possibly after the dangling-node repair of Equation~\eqref{eq:transition}. For $d\in(0,1)$ the Google matrix of Equation~\eqref{eq:google-matrix},
\[
G_{\mathrm{PR}}=d P^{\top}+\frac{1-d}{|V|}\mathbf{1}\mathbf{1}^{\top},
\]
is positive and column-stochastic. By the Perron--Frobenius theorem for positive matrices, $G_{\mathrm{PR}}$ has a single eigenvalue of modulus one, and the corresponding eigenvector $\mathbf{r}$ is strictly positive; we can scale it so that $\sum_i r_i=1$. It follows that, once the damping and the dangling-node rules are fixed, every directed graph has unique PageRank scores \parencite{langville2006pagerank}.

Three restart vectors in this thesis are not uniform. AWP's, proportional to activeness, is zero at every contract that has sent nothing, that of EndorseRank with activity restarts is zero at every contract that has approved nothing, and S-PR's, the normalized score of the allowance walk, is zero at every node of the transfer layer that has no allowance score, so none of the three Google matrices is positive. The fixed point is still unique: for two probability vectors the restart term cancels in their difference, and the operator multiplies their $\|\cdot\|_1$ distance by at most $d$. In general it is not strictly positive, however, because a node with no restart mass that no path from a restart node reaches keeps a score of zero. Uniqueness means that, once its parameters are fixed, each method assigns exactly one score vector to a given edge set, so any difference between $\mathbf{s}_{\mathrm{ER}}$ and $\mathbf{s}_{\mathrm{AWP}}$ arises from their edges and restart distributions.

For fixed $d$ the scores also depend continuously on the entries of $P$: a second row-stochastic matrix $\tilde P$ on the same nodes moves them by at most $\frac{d}{1-d}\max_i\|P_{i\cdot}-\tilde P_{i\cdot}\|_1$ in the $\|\cdot\|_1$ norm \parencite{langville2006pagerank}. A row of $P$ is not continuous in the edge weights, however: it jumps between the uniform dangling row and a row concentrated on its out-neighbors when a node gains its first out-edge or loses its last one, however small the weight involved.\label{sec:pr-continuity}

\dissertationSubheading[sec:pr-sparsity]{Sparsity and per-iteration cost}

With $m=|E|$ nonzero edges and $n=|V|$ nodes, the sparse product $P^{\top}\mathbf{x}$ costs $O(m)$ arithmetic operations, each iteration also updates vectors of length $n$, and building the sparse operator from the edge list is $O(m)$, so power iteration over $I$ iterations is $O(m+(m+n)I)$. Both scores run the same solver. In the measured calls an iteration of AWP takes about $15.9$ times as long as one of EndorseRank, less than the edge ratio of $22.4$, and AWP needs more iterations (Section~\ref{sec:benchmark-results}), so neither $m$ nor $I$ alone sets the time ratio.

\dissertationSubheading[sec:kendall-properties]{Properties of Spearman's \texorpdfstring{$\rho$}{rho} and Kendall's \texorpdfstring{$\tau$}{tau}}
\label{sec:spearman-properties}

Spearman's $\rho$ and Kendall's $\tau_b$ lie in $[-1,1]$, and the evaluation pipeline computes both with standard tie-aware implementations. Across the fourteen pairs of EndorseRank or AWP with a proxy (Tables~\ref{tab:alignment-transfer} to~\ref{tab:alignment-sybil-stability}) both coefficients are positive and $\rho$ is the larger of the two. The denominator of $\tau_b$ in Equation~\eqref{eq:kendall} differs from $n_0$ only through ties, and among the scores of this thesis ties are frequent: a contract outside the graph of a score scores zero, so on the matched cohort EndorseRank ties $65.4\%$ of the contract pairs and AWP $13.7\%$ (Section~\ref{sec:rank-divergence}), and the integer proxies tie many pairs as well. A contrast between two scores on one proxy therefore also reflects which contracts each score leaves at zero. Table~\ref{tab:tie-sensitivity} repeats the family means with the contracts outside a graph kept as isolated nodes, and Table~\ref{tab:union-members} the main out-of-window contrasts on the members inside the union graph.

% !TEX root = ../main.tex
\section{Data Description}
\label{ch:appendix-data}

This appendix describes the data behind Chapter~\ref{ch:results}: the events and the layers of the contract graph, the five tokens with their prices and allowance amounts, the account types, balances and labels of the contract graph, and the detail of the discovery-order test.

\dissertationSubheading[sec:data-events]{Extracted events and graphs}

The events of the main analysis are the \texttt{Approval} logs of any value and the \texttt{Transfer} logs of non-zero value that the five tokens emitted from 1~April~2025 to 30~September~2026, the months that the two score windows and the two label windows cover; Section~\ref{sec:dataset} gives their number. They form $73{,}058{,}133$ distinct ordered pairs of addresses. Only the rows with a contract at both ends enter a graph (Table~\ref{tab:graph-layers}). The node table numbers $1{,}403{,}959$ contracts over both freezes: every contract that ends a pair of either layer, with or without the balance cap, and every contract of the cohort and spender lists. A further $152{,}257{,}072$ rows of the two score windows lead from a user account to a contract or to an address that was not typed; no score reads them, and the importance measures of the discovery-order test are computed from them (Section~\ref{sec:discovery-protocol}). The registered analysis extracted the events of every ERC-20 token with a contract of the plan's cohort at one end, from 1~October~2023 to 30~June~2026: $389{,}866{,}882$ approval logs, zero-value ones included, and $1{,}538{,}404{,}067$ transfer logs of non-zero value.

\dissertationSubheading[sec:data-tokens]{Tokens, prices and allowance amounts}

Section~\ref{sec:amount-revision} gives the rule by which the five tokens were chosen. Table~\ref{tab:usd-tokens} gives for each of them the address, the decimals and the Chainlink feed, the positive latest allowances between contracts at $T_{\mathrm{obs}}$ before and after the balance cap, and the data behind the daily prices.

% Generated by experiments/scripts/export_latex.py (export_latex_c12.py) -- do not edit by hand.
% Source: experiments/data/0-usd-token-list-and-prices/listed_tokens.csv
% Source: experiments/data/2-processed-tables-and-evaluations/revised-contract-12m/describe/describe.json
% Source: experiments/data/0-usd-token-list-and-prices/chainlink_check.json
\begin{table}[htbp]
\centering
\caption[The five tokens, their allowances between contracts and their prices.]{The five tokens of the analysis, with their Chainlink feed (market-risk category), the positive latest allowances between contracts at $T_{\mathrm{obs}}$ and those that stay positive after the balance cap, the days of the price series filled from the second DefiLlama key, filled from the previous day and replaced as one-day errors (of the 548 days from 1~April~2025 to 30~September~2026 that the windows cover), and the median and largest absolute difference, in percent, between the DefiLlama price and the Chainlink feed read on the chain on the 20 check dates among those days. The five-token events of the typing window number 1{,}240{,}275{,}042; 603{,}475{,}784 of them have a contract at both ends.}
\label{tab:usd-tokens}
\fitwidth{%
\begin{tabular}{llrlrrrrrr}
\toprule
Token & Address & Dec. & Chainlink feed & Allowances & After the cap & 2nd key & Filled & Errors & CL diff.\ (\%) \\
\midrule
USDC & \texttt{0xaf88\ldots{}5831} & 6 & USDC / USD (low) & 255{,}339 & 108{,}916 & 0 & 0 & 0 & 0.01 / 0.02 \\
USDT0 & \texttt{0xfd08\ldots{}cbb9} & 6 & USDT / USD (low) & 36{,}468 & 5{,}648 & 0 & 0 & 0 & 0.01 / 0.15 \\
WETH & \texttt{0x82af\ldots{}bab1} & 18 & ETH / USD (low) & 30{,}826 & 1{,}395 & 0 & 0 & 0 & 0.04 / 0.14 \\
WBTC & \texttt{0x2f2a\ldots{}5b0f} & 8 & WBTC / USD (low) & 11{,}841 & 386 & 0 & 0 & 0 & 0.05 / 0.24 \\
ARB & \texttt{0x912c\ldots{}6548} & 18 & ARB / USD (low) & 4{,}808 & 378 & 0 & 0 & 0 & 0.06 / 0.24 \\
\bottomrule
\end{tabular}
}
\end{table}


Each of the five tokens has a DefiLlama price under its own Arbitrum One address on each of the $548$ days from 1~April~2025 to 30~September~2026: no day is taken from a second key of the service, filled from the previous day or corrected as a one-day error. Against the Chainlink feeds read on the chain on $20$ of these days per token, the DefiLlama prices differ by a median of $0.01\%$ to $0.06\%$ per token and by at most $0.24\%$.

Unlimited allowances are common. Of the $339{,}282$ positive latest allowances between contracts at $T_{\mathrm{obs}}$, $223{,}892$ ($66.0\%$) are of $1.15\times10^{77}$ base units or more, close to $2^{256}-1$, and $230{,}313$, the unlimited ones included, are worth at least the cap $U$ before the balance cap (Table~\ref{tab:allowance-cap} gives the unlimited allowances per token). In the registered analysis, which kept every token in raw base units, $99.52\%$ of the ego transfers and $99.78\%$ of the $28{,}069{,}418$ latest allowances at $T_{\mathrm{obs}}$ were of ten base units or more, where $V$ with $b=1$ is within $10^{-4}$ of one, and $1{,}532{,}301$ of the $13{,}029{,}306$ owners with a positive latest allowance held both an unlimited and a smaller allowance.

\dissertationSubheading[sec:data-contract-graph]{Contract graph: account types, smart accounts and balances}

Table~\ref{tab:contract-types} gives the account types behind the contract graph; Section~\ref{sub:account-types} states the rule and Appendix~\ref{sec:account-type-classification} the reads.

% Generated by experiments/scripts/export_latex.py (export_latex_c12.py) -- do not edit by hand.
% Source: experiments/data/2-processed-tables-and-evaluations/revised-contract-12m/describe/describe.json
\begin{table}[htbp]
\centering
\caption[Account types of the contract graph.]{Account types of the contract graph. The five-token events of the typing window have 25{,}537{,}280 distinct addresses other than the zero address; 17{,}477{,}779 of them sent a transaction and are user accounts by that fact. The table types the 6{,}230{,}270 addresses that were read with \texttt{eth\_getCode} at block 513{,}216{,}506: every address that sent no transaction and has a counterparty that is neither a transaction sender nor the zero address, the members of the cohorts and two audit samples. An address outside this set that sent no transaction has only user accounts as counterparties and cannot end an edge between contracts. Smart accounts are user accounts: a Safe answers its threshold and owners, and an ERC-4337 account answers a canonical EntryPoint and carries a \texttt{validateUserOp} selector and no paymaster selector; 192{,}656 further accounts whose proxy hides that selector were moved by listed code families (37{,}337 by code hash and 155{,}319 by code hash and implementation). Of the addresses read, 524{,}454 are ERC-1967 proxies, 503{,}516 EIP-1167 clones and 45{,}362 Safe proxies. With the transaction senders, 21{,}877{,}279 addresses are typed user accounts and 1{,}829{,}719 contracts.}
\label{tab:contract-types}
\fitwidth{%
\begin{tabular}{llr}
\toprule
Code read at the block & Class & Addresses \\
\midrule
Externally owned account without code & user account & 2{,}909{,}509 \\
Externally owned account with an EIP-7702 delegation & user account & 1{,}165{,}025 \\
Safe (threshold and owners answered) & user account & 46{,}204 \\
ERC-4337 account, by the rule of the reads & user account & 87{,}157 \\
ERC-4337 account, by the listed code families & user account & 192{,}656 \\
Any other code & contract (graph node) & 1{,}829{,}719 \\
\midrule
Addresses read & & 6{,}230{,}270 \\
\bottomrule
\end{tabular}
}
\end{table}


The reads alone miss ERC-4337 accounts whose proxy hides the \texttt{validateUserOp} selector. Two lists of code families, drawn from the code alone without reading a label, a cohort membership or a score of the addresses they move, type $192{,}656$ more addresses as smart accounts. The first holds $20$ code hashes of small proxies whose implementation the reads did not find and whose representative answers a canonical EntryPoint ($37{,}337$ addresses). The second holds seven pairs of a code hash and an implementation of ERC-1967 proxies whose implementation carries the selector and whose getter \texttt{entrypoint()} answers a canonical EntryPoint where \texttt{entryPoint()} answers none ($155{,}319$ addresses).

Kinds of smart account that the rule does not find stay contracts and graph nodes. After the two lists, $115{,}084$ typed contracts still carry a \texttt{validateUserOp} selector and no paymaster selector in the code examined: $104{,}187$ of them share one code hash whose getter \texttt{entrypoint()} reverts or does not exist, and $8{,}793$ run one implementation whose EntryPoint is not a canonical one. Contracts whose code suggests a smart account of another kind, neither a Safe nor an ERC-4337 account, stay as well, as do the contracts that services create for each payment, order or user, which no rule picks out.

The balance of an owner in a token at a freeze is the net flow of the token's \texttt{Transfer} logs to the owner from the first block of the chain up to the freeze, set to zero when it is negative or below $10^{-6}$ token (Section~\ref{sub:balances}). It was computed for the $6{,}033{,}263$ pairs of a token and an owner with an approval of that token from 1~April~2025 to $T_{\mathrm{obs}}$. A check compared $1{,}500$ of these pairs, for each token $150$ with a zero and $150$ with a positive reconstructed balance at $T_{\mathrm{obs}}$, with \texttt{balanceOf} at block $479{,}089{,}705$, the last block of 30~June~2026: $1{,}493$ agree ($99.5\%$, against a threshold of $98\%$). For the other seven, four of WETH and three of USDT0, each with a reconstructed balance that counts as zero, the call returned no value, and the check counts them as disagreements; no pair with an answer differs.

Table~\ref{tab:allowance-cap} gives the effect of the balance cap on the latest allowances between contracts, per token. Most of the allowances that stay positive are counted at a balance below the allowed amount ($114{,}454$ of $116{,}723$ at $T_{\mathrm{obs}}$), and the two freezes differ in USDC, of which $108{,}916$ allowances stay positive at $T_{\mathrm{obs}}$ and $11{,}381$ at $t_1$. After the cap $150$ allowances at $T_{\mathrm{obs}}$ and $204$ at $t_1$ are worth $U$ or more, so the cap $U$ of the allowance layer of C-PR binds on few of them. The labels, unlike the graphs, count every account, and Table~\ref{tab:label-split} splits each label of the two label windows by the class of the account it counts (Section~\ref{sec:holdout-results}).

% Generated by experiments/scripts/export_latex.py (export_latex_c12.py) -- do not edit by hand.
% Source: experiments/data/2-processed-tables-and-evaluations/revised-contract-12m/describe/describe.json
\begin{table}[htbp]
\centering
\caption[Latest allowances between contracts and the balance cap.]{Latest allowances between contracts and the balance cap, per token and freeze. An allowance is the latest \texttt{Approval} of a (token, owner, spender) triple in the score window, kept when its value is positive. It is counted up to the owner's balance of the token at the freeze: unlimited allowances are those of at least $1.1\times10^{77}$ base units, an allowance is capped to zero when the owner holds no balance of the token at the freeze, and capped below the allowance when the balance is positive and smaller than the allowance. Only the allowances that stay positive enter the allowance layer. The last column is the sum of the capped values in US dollars. At $T_{\mathrm{obs}}$ 65.6\% of the allowances are capped to zero and 66.0\% are unlimited (at $t_1$: 93.6\% and 67.4\%). The median positive capped value at $t_1$ is \$0.071, the constant $m_A$.}
\label{tab:allowance-cap}
\fitwidth{%
\begin{tabular}{lrrrrrr}
\toprule
Token & Positive allowances & Unlimited & Capped to zero & Capped below & Positive after the cap & Capped value (USD) \\
\midrule
\multicolumn{7}{l}{\emph{At $T_{\mathrm{obs}}$}} \\
USDC & 255{,}339 & 167{,}068 & 146{,}423 & 106{,}988 & 108{,}916 & $3.3\times10^{7}$ \\
USDT0 & 36{,}468 & 29{,}609 & 30{,}820 & 5{,}522 & 5{,}648 & $6.5\times10^{6}$ \\
WETH & 30{,}826 & 17{,}608 & 29{,}431 & 1{,}270 & 1{,}395 & $2.0\times10^{7}$ \\
WBTC & 11{,}841 & 5{,}592 & 11{,}455 & 341 & 386 & $7.8\times10^{6}$ \\
ARB & 4{,}808 & 4{,}015 & 4{,}430 & 333 & 378 & $8.4\times10^{5}$ \\
All five tokens & 339{,}282 & 223{,}892 & 222{,}559 & 114{,}454 & 116{,}723 & $6.8\times10^{7}$ \\
\midrule
\multicolumn{7}{l}{\emph{At $t_1$}} \\
USDC & 223{,}685 & 146{,}928 & 212{,}304 & 9{,}640 & 11{,}381 & $3.3\times10^{7}$ \\
WETH & 36{,}916 & 24{,}198 & 35{,}474 & 1{,}294 & 1{,}442 & $3.2\times10^{7}$ \\
USDT0 & 33{,}673 & 27{,}950 & 27{,}666 & 5{,}851 & 6{,}007 & $7.6\times10^{6}$ \\
WBTC & 9{,}853 & 5{,}409 & 9{,}413 & 390 & 440 & $1.1\times10^{7}$ \\
ARB & 4{,}795 & 3{,}760 & 4{,}442 & 305 & 353 & $6.1\times10^{5}$ \\
All five tokens & 308{,}922 & 208{,}245 & 289{,}299 & 17{,}480 & 19{,}623 & $8.4\times10^{7}$ \\
\bottomrule
\end{tabular}
}
\end{table}


% Generated by experiments/scripts/export_latex.py (export_latex_c12.py) -- do not edit by hand.
% Source: experiments/data/2-processed-tables-and-evaluations/revised-contract-12m/holdout-w0/eval_summary.json
% Source: experiments/data/2-processed-tables-and-evaluations/revised-contract-12m/window-w1/eval_summary.json
\begin{table}[htbp]
\centering
\caption[The labels by the class of the account they count.]{The labels of the two label windows by the class of the account they count: for each label, its total over the spender cohort, the parts counted on user accounts and on contracts, the share of the contracts, and the spenders that have a positive part of each kind. An approval pair is counted by the class of its owner and a sender by its own class; a handful of senders are neither (the zero address of mints), so the two parts can fall short of the total by them. The split is a description outside every rule: the labels are those of the twelve-month analysis, unchanged, and most of what they count are user accounts, which no score of the contract graph reads.}
\label{tab:label-split}
\fitwidth{%
\begin{tabular}{lrrrrrr}
\toprule
Label & Total & From user accounts & From contracts & Contracts' share & Spenders with user accounts & Spenders with contracts \\
\midrule
\multicolumn{7}{l}{\emph{W0 ($n=7{,}864$ spender contracts)}} \\
New approval pairs & 748{,}258 & 676{,}525 & 71{,}733 & 9.6\% & 1{,}322 & 876 \\
New transfer senders & 651{,}474 & 562{,}007 & 89{,}458 & 13.7\% & 1{,}044 & 883 \\
Revocations & 61{,}869 & 60{,}111 & 1{,}758 & 2.8\% & 947 & 287 \\
Drained owners & 570{,}724 & 547{,}427 & 23{,}297 & 4.1\% & 1{,}271 & 926 \\
\midrule
\multicolumn{7}{l}{\emph{W1 ($n=7{,}915$ spender contracts)}} \\
New approval pairs & 702{,}275 & 617{,}088 & 85{,}187 & 12.1\% & 1{,}220 & 641 \\
New transfer senders & 618{,}461 & 509{,}194 & 109{,}258 & 17.7\% & 970 & 855 \\
Revocations & 43{,}625 & 42{,}270 & 1{,}355 & 3.1\% & 761 & 279 \\
Drained owners & 528{,}557 & 503{,}788 & 24{,}769 & 4.7\% & 1{,}139 & 847 \\
\bottomrule
\end{tabular}
}
\end{table}


\dissertationSubheading[sec:discovery-detail]{Discovery-order test: universes, coverage, sensitivity and diagnostics}

This section holds the tables of the discovery-order test that Section~\ref{sec:discovery-results} reads; Section~\ref{sec:discovery-protocol} defines their terms. The test was defined after the registered results were known, and it decides no rule of the analysis plan or the registration.

Table~\ref{tab:discovery-sizes} gives the sizes of the contract graph, of its bow-tie and of the two universes at both freezes, with the sizes of the hot sets. Table~\ref{tab:discovery-coverage} shows how much of the two importances each universe holds: at $T_{\mathrm{obs}}$ the universe of the \texttt{both} mode holds $96.5\%$ of $I_T$ and $93.0\%$ of $I_A$ by value, but $64.2\%$ and $38.1\%$ of the contracts that make up the top $1\%$ by each.

% Generated by experiments/scripts/export_latex.py (export_latex_c12.py) -- do not edit by hand.
% Source: experiments/data/2-processed-tables-and-evaluations/revised-contract-12m/discovery/discovery_tobs_both.json
% Source: experiments/data/2-processed-tables-and-evaluations/revised-contract-12m/discovery/discovery_tobs_out.json
% Source: experiments/data/2-processed-tables-and-evaluations/revised-contract-12m/discovery/discovery_t1_both.json
% Source: experiments/data/2-processed-tables-and-evaluations/revised-contract-12m/discovery/discovery_t1_out.json
\begin{table}[htbp]
\centering
\caption[Sizes of the contract graph and of the universes of the discovery-order test.]{Sizes of the contract graph and of the two universes of the discovery-order test at both freezes. The union graph holds the distinct ordered pairs of different contracts with an edge in either layer; the bow-tie is that of the union graph. $\mathrm{SCC}\cup\mathrm{OUT}$ holds fewer than half the contracts of the weakly connected component of the giant SCC (WCC) at $t_1$, so by the rule of the protocol the \texttt{both} mode, whose universe is the WCC, is the primary mode at both freezes and the \texttt{out} mode the sensitivity. A hot set $H_q(I)$ holds the contracts of the universe with the largest positive importance $I$, $q=1/100$ of the universe with ties included, so it is smaller when fewer contracts have a positive importance.}
\label{tab:discovery-sizes}
\fitwidth{%
\begin{tabular}{lrr}
\toprule
Quantity & $T_{\mathrm{obs}}$ (primary) & $t_1$ \\
\midrule
Contracts of $G$ (one numbering for both freezes) & 1{,}403{,}959 & 1{,}403{,}959 \\
Capped allowance pairs & 115{,}006 & 17{,}940 \\
Transfer pairs & 2{,}580{,}709 & 2{,}269{,}073 \\
Union graph: contracts with a pair & 1{,}299{,}438 & 1{,}111{,}554 \\
Union graph: pairs & 2{,}575{,}632 & 2{,}269{,}421 \\
\midrule
Giant strongly connected component (SCC) & 265{,}804 & 222{,}592 \\
IN (reaches the SCC) & 174{,}239 & 169{,}237 \\
OUT (reached from the SCC) & 242{,}833 & 109{,}804 \\
Other contracts with a pair & 616{,}562 & 609{,}921 \\
\midrule
$|\mathrm{SCC}\cup\mathrm{OUT}|$, universe of the \texttt{out} mode & 508{,}637 & 332{,}396 \\
Union pairs inside it & 1{,}546{,}906 & 1{,}283{,}946 \\
$|\mathrm{WCC}|$, universe of the \texttt{both} mode (primary) & 1{,}169{,}044 & 1{,}015{,}867 \\
Union pairs inside it & 2{,}378{,}418 & 2{,}121{,}895 \\
Capped allowance pairs inside the WCC & 16{,}776 & 17{,}741 \\
\midrule
$|H_q(I_T)|$, \texttt{both} mode & 13{,}141 & 12{,}697 \\
$|H_q(I_A)|$, \texttt{both} mode & 11{,}662 & 8{,}136 \\
$|H_q(I_T)|$, \texttt{out} mode & 5{,}087 & 3{,}324 \\
$|H_q(I_A)|$, \texttt{out} mode & 5{,}088 & 3{,}324 \\
\bottomrule
\end{tabular}
}
\end{table}


% Generated by experiments/scripts/export_latex.py (export_latex_c12.py) -- do not edit by hand.
% Source: experiments/data/2-processed-tables-and-evaluations/revised-contract-12m/discovery/discovery_tobs_both.json
% Source: experiments/data/2-processed-tables-and-evaluations/revised-contract-12m/discovery/discovery_tobs_out.json
% Source: experiments/data/2-processed-tables-and-evaluations/revised-contract-12m/discovery/discovery_t1_both.json
% Source: experiments/data/2-processed-tables-and-evaluations/revised-contract-12m/discovery/discovery_t1_out.json
\begin{table}[htbp]
\centering
\caption[Coverage of the two importances by the contract graph and the universes.]{Coverage of the two importances by the contract graph and by the universes of the discovery-order test. $I_T$ is the US-dollar value that user accounts transferred to a contract in the score window and $I_A$ the value they had authorized it to spend at the freeze, each counted up to $U_{\mathrm{user}}$ per account and, for $I_A$, up to the account's balance. At $T_{\mathrm{obs}}$ 809{,}296 typed contracts have a positive importance, 112{,}210 of them outside $G$ (at $t_1$: 715{,}480 and 109{,}094). The parts of the bow-tie are those of the union graph and sum to 100\%. User accounts also transferred \$1.68 billion at $T_{\mathrm{obs}}$ and \$1.37 billion at $t_1$ to addresses that were not typed because every counterparty of theirs is a user account; these addresses are never graph nodes, and the sum bounds the inflow that the coverage cannot place.}
\label{tab:discovery-coverage}
\fitwidth{%
\begin{tabular}{lrrrr}
\toprule
 & \multicolumn{2}{c}{$T_{\mathrm{obs}}$} & \multicolumn{2}{c}{$t_1$} \\
\cmidrule(lr){2-3} \cmidrule(lr){4-5}
Share of the importance & $I_T$ & $I_A$ & $I_T$ & $I_A$ \\
\midrule
Sum over all typed contracts & \$16.61 billion & \$309.5 million & \$16.62 billion & \$343.8 million \\
\midrule
Giant SCC & 71.2\% & 50.2\% & 74.1\% & 49.1\% \\
IN & 9.7\% & 1.7\% & 9.0\% & 2.6\% \\
OUT & 10.1\% & 40.4\% & 6.6\% & 38.2\% \\
Other contracts with a pair & 8.1\% & 3.4\% & 9.5\% & 6.3\% \\
Contracts of $G$ without a union pair & 0.4\% & 3.6\% & 0.5\% & 3.3\% \\
Contracts outside $G$ & 0.4\% & 0.7\% & 0.3\% & 0.5\% \\
\midrule
Held by the WCC (\texttt{both} mode) & 96.5\% & 93.0\% & 97.0\% & 93.8\% \\
Held by $\mathrm{SCC}\cup\mathrm{OUT}$ (\texttt{out} mode) & 81.3\% & 90.6\% & 80.7\% & 87.3\% \\
Top 1\% of the contracts by the importance: share in the WCC & 64.2\% & 38.1\% & 64.8\% & 34.6\% \\
Top 1\%: share in $\mathrm{SCC}\cup\mathrm{OUT}$ & 29.6\% & 26.2\% & 32.5\% & 23.7\% \\
\bottomrule
\end{tabular}
}
\end{table}


Table~\ref{tab:discovery-sensitivity} changes one setting of the deciding run at a time. In no row does the interval of $D_2$ lie above zero; the interval of $D_1$ on $H_q(I_A)$ lies below zero in two rows and contains zero in two others. Table~\ref{tab:discovery-diagnostics} gives the share of picks that each ordering decides by a tie, the share of early visits that go to contracts without user-account demand and the visit at which each ordering reaches the hubs: the median hub is visit $950$ under AWP$'$ and $4{,}088$ under C-PR$'$, against $25{,}565$ under $\mathrm{IB}'_U$.

% Generated by experiments/scripts/export_latex.py (export_latex_c12.py) -- do not edit by hand.
% Source: experiments/data/2-processed-tables-and-evaluations/revised-contract-12m/discovery/discovery_tobs_both.json
\begin{table}[htbp]
\centering
\caption[Discovery-order test: sensitivity of the decision contrasts.]{Discovery-order test at $T_{\mathrm{obs}}$ in the primary mode: sensitivity of the two decision contrasts, $D_1$ (C-PR$'$ minus BFS) and $D_2$ (C-PR$'$ minus $\mathrm{IB}'_U$), each over its 95\% percentile intervals of a paired bootstrap over the hot contracts (400 resamples, seed 42). The rows change one setting at a time: the share $x$ of the universe over which the area is taken, the share $q$ of the hot set, the hot set without the hubs (the contracts with the highest union in-degree, removed from the hot set but kept in the graph), the importance measured by the number of user accounts with at least one dollar, a bootstrap over the seeds in place of the hot contracts, and a schedule that recomputes the ordering four times as often, on the first seeds. Only 11{,}662 contracts of the universe have a positive $I_A$, so the hot set of $I_A$ is the same at both shares $q$. The other mode and the other freeze are in Table~\ref{tab:discovery-decision}. None of these rows carries a verdict. The test was defined after the registered results were known; it decides no rule of the analysis plan or the registration.}
\label{tab:discovery-sensitivity}
\fitwidth{%
\begin{tabular}{lcccc}
\toprule
 & \multicolumn{2}{c}{On $H_q(I_T)$} & \multicolumn{2}{c}{On $H_q(I_A)$} \\
\cmidrule(lr){2-3} \cmidrule(lr){4-5}
Setting & $D_1$ & $D_2$ & $D_1$ & $D_2$ \\
\midrule
Decision ($x=10\%$, $q=1\%$) & \begin{tabular}[c]{@{}c@{}}+0.0520\\[-1pt]{\footnotesize [+0.0485, +0.0558]}\end{tabular} & \begin{tabular}[c]{@{}c@{}}-0.0051\\[-1pt]{\footnotesize [-0.0087, -0.0010]}\end{tabular} & \begin{tabular}[c]{@{}c@{}}+0.0061\\[-1pt]{\footnotesize [+0.0014, +0.0111]}\end{tabular} & \begin{tabular}[c]{@{}c@{}}-0.0571\\[-1pt]{\footnotesize [-0.0621, -0.0519]}\end{tabular} \\
\midrule
$x=5\%$ & \begin{tabular}[c]{@{}c@{}}+0.0673\\[-1pt]{\footnotesize [+0.0635, +0.0714]}\end{tabular} & \begin{tabular}[c]{@{}c@{}}+0.0026\\[-1pt]{\footnotesize [-0.0010, +0.0064]}\end{tabular} & \begin{tabular}[c]{@{}c@{}}+0.0472\\[-1pt]{\footnotesize [+0.0429, +0.0521]}\end{tabular} & \begin{tabular}[c]{@{}c@{}}-0.0132\\[-1pt]{\footnotesize [-0.0178, -0.0084]}\end{tabular} \\
$x=20\%$ & \begin{tabular}[c]{@{}c@{}}+0.0199\\[-1pt]{\footnotesize [+0.0169, +0.0230]}\end{tabular} & \begin{tabular}[c]{@{}c@{}}-0.0063\\[-1pt]{\footnotesize [-0.0091, -0.0034]}\end{tabular} & \begin{tabular}[c]{@{}c@{}}-0.0518\\[-1pt]{\footnotesize [-0.0575, -0.0460]}\end{tabular} & \begin{tabular}[c]{@{}c@{}}-0.0966\\[-1pt]{\footnotesize [-0.1023, -0.0905]}\end{tabular} \\
$q=5\%$ & \begin{tabular}[c]{@{}c@{}}+0.0132\\[-1pt]{\footnotesize [+0.0121, +0.0144]}\end{tabular} & \begin{tabular}[c]{@{}c@{}}-0.0067\\[-1pt]{\footnotesize [-0.0078, -0.0057]}\end{tabular} & \begin{tabular}[c]{@{}c@{}}+0.0061\\[-1pt]{\footnotesize [+0.0014, +0.0111]}\end{tabular} & \begin{tabular}[c]{@{}c@{}}-0.0571\\[-1pt]{\footnotesize [-0.0621, -0.0519]}\end{tabular} \\
Without the 1{,}170 hubs & \begin{tabular}[c]{@{}c@{}}+0.0306\\[-1pt]{\footnotesize [+0.0274, +0.0341]}\end{tabular} & \begin{tabular}[c]{@{}c@{}}-0.0126\\[-1pt]{\footnotesize [-0.0159, -0.0088]}\end{tabular} & \begin{tabular}[c]{@{}c@{}}-0.0019\\[-1pt]{\footnotesize [-0.0065, +0.0026]}\end{tabular} & \begin{tabular}[c]{@{}c@{}}-0.0607\\[-1pt]{\footnotesize [-0.0654, -0.0553]}\end{tabular} \\
Importance by counts ($N_T$, $N_A$) & \begin{tabular}[c]{@{}c@{}}+0.0188\\[-1pt]{\footnotesize [+0.0163, +0.0211]}\end{tabular} & \begin{tabular}[c]{@{}c@{}}-0.0282\\[-1pt]{\footnotesize [-0.0308, -0.0254]}\end{tabular} & \begin{tabular}[c]{@{}c@{}}+0.0191\\[-1pt]{\footnotesize [+0.0129, +0.0262]}\end{tabular} & \begin{tabular}[c]{@{}c@{}}-0.0473\\[-1pt]{\footnotesize [-0.0537, -0.0408]}\end{tabular} \\
Bootstrap over seeds & \begin{tabular}[c]{@{}c@{}}+0.0520\\[-1pt]{\footnotesize [+0.0320, +0.0717]}\end{tabular} & \begin{tabular}[c]{@{}c@{}}-0.0051\\[-1pt]{\footnotesize [-0.0249, +0.0179]}\end{tabular} & \begin{tabular}[c]{@{}c@{}}+0.0061\\[-1pt]{\footnotesize [-0.0315, +0.0398]}\end{tabular} & \begin{tabular}[c]{@{}c@{}}-0.0571\\[-1pt]{\footnotesize [-0.0873, -0.0221]}\end{tabular} \\
$\rho/4$ (recompute after 1/200 more visits; 5 seeds) & \begin{tabular}[c]{@{}c@{}}+0.0381\\[-1pt]{\footnotesize [+0.0347, +0.0418]}\end{tabular} & \begin{tabular}[c]{@{}c@{}}-0.0129\\[-1pt]{\footnotesize [-0.0167, -0.0087]}\end{tabular} & \begin{tabular}[c]{@{}c@{}}-0.0134\\[-1pt]{\footnotesize [-0.0181, -0.0082]}\end{tabular} & \begin{tabular}[c]{@{}c@{}}-0.0713\\[-1pt]{\footnotesize [-0.0765, -0.0658]}\end{tabular} \\
\bottomrule
\end{tabular}
}
\end{table}


% Generated by experiments/scripts/export_latex.py (export_latex_c12.py) -- do not edit by hand.
% Source: experiments/data/2-processed-tables-and-evaluations/revised-contract-12m/discovery/discovery_tobs_both.json
% Source: experiments/data/2-processed-tables-and-evaluations/revised-contract-12m/discovery/discovery_tobs_out.json
\begin{table}[htbp]
\centering
\caption[Discovery-order test: ties, wasted visits and hub visits.]{Discovery-order test at $T_{\mathrm{obs}}$: diagnostics of the crawls. Ties: the share of picks whose group of tied scores held more than one contract, in the primary mode (\texttt{both}) and in the \texttt{out} mode; a tie is broken by the earlier discovery and then by the lower address. Wasted visits: the share of the first 1\% and of the first 10\% of visits that went to contracts with $I_T = I_A = 0$ (primary mode). Hub visit: over the 1{,}170 hubs of the primary universe, the median of each hub's mean visit index across seeds; the largest hub has a union in-degree of 441{,}612. The count orderings and the walks on the allowance layer alone decide nearly every pick by a tie.}
\label{tab:discovery-diagnostics}
\fitwidth{%
\begin{tabular}{lccccr}
\toprule
 & \multicolumn{2}{c}{Picks decided by a tie} & \multicolumn{2}{c}{Wasted visits} & Median hub \\
\cmidrule(lr){2-3} \cmidrule(lr){4-5}
Ordering & \texttt{both} & \texttt{out} & first 1\% & first 10\% & visit index \\
\midrule
Breadth-first (BFS) & 0.000 & 0.000 & 0.891 & 0.845 & 73{,}219 \\
Random & 0.000 & 0.000 & 0.841 & 0.682 & 272{,}806 \\
\midrule
$\mathrm{IB}'_A$ (allowance backlinks) & 1.000 & 1.000 & 0.874 & 0.862 & 59{,}089 \\
$\mathrm{IB}'_T$ (transfer backlinks) & 1.000 & 0.999 & 0.678 & 0.888 & 25{,}566 \\
$\mathrm{IB}'_U$ (contract backlinks) & 1.000 & 0.999 & 0.675 & 0.888 & 25{,}565 \\
\midrule
EndorseRank$'$ & 0.999 & 0.998 & 0.874 & 0.863 & 60{,}494 \\
EndorseRank-activity$'$ & 0.991 & 0.997 & 0.692 & 0.857 & 66{,}699 \\
AWP$'$ & 0.252 & 0.005 & 0.522 & 0.835 & 950 \\
\midrule
C-PR$'$ ($\lambda=1$) & 1.000 & 0.999 & 0.874 & 0.863 & 60{,}840 \\
C-PR$'$ ($\lambda=0$) & 0.531 & 0.019 & 0.747 & 0.911 & 5{,}156 \\
C-PR$'$ ($\lambda=0.25$) & 0.534 & 0.019 & 0.727 & 0.908 & 4{,}185 \\
C-PR$'$ ($\lambda=0.5$) & 0.535 & 0.023 & 0.719 & 0.900 & 4{,}088 \\
C-PR$'$ ($\lambda=0.75$) & 0.537 & 0.025 & 0.719 & 0.889 & 4{,}203 \\
S-PR$'$ & 0.531 & 0.019 & 0.743 & 0.910 & 4{,}885 \\
\midrule
$\delta'$ (normalized in-strength) & 0.634 & 0.283 & 0.497 & 0.650 & 1{,}404 \\
\bottomrule
\end{tabular}
}
\end{table}


Figure~\ref{fig:discovery-curves-out} shows the curves of the \texttt{out} mode at $T_{\mathrm{obs}}$, which carries no verdict. In that mode C-PR$'$ reaches the hot contracts sooner than the breadth-first crawl and than $\mathrm{IB}'_U$ on both importances (Table~\ref{tab:discovery-decision}), but it does not lead the walks on the transfer layer alone: on $H_q(I_T)$ it lies below AWP$'$ ($-0.0617$, $[-0.0709,-0.0526]$) and below its own transfer layer, C-PR$'$ ($\lambda=0$) ($-0.0305$, $[-0.0343,-0.0264]$), and on $H_q(I_A)$ it is level with AWP$'$ ($-0.0016$, $[-0.0091,+0.0063]$) and above C-PR$'$ ($\lambda=0$) ($+0.0153$, $[+0.0108,+0.0206]$).

\begin{figure}[htbp]
\centering
\includegraphics[width=\textwidth]{\figpath{discovery-curves-out.pdf}}
\caption[{Discovery-order test: share of the hot contracts visited when a visit reveals out-edges only.}]{Discovery-order test at $T_{\mathrm{obs}}$ in the \texttt{out} mode, which carries no verdict: the share of the hot contracts visited, $P_{ST}$, against the contracts visited as a percentage of the universe, for the hot sets of user-account inflow $I_T$ (left) and of user-authorised value $I_A$ (right). In this mode a visit reveals out-edges only, and the universe is the giant strongly connected component with the contracts reachable from it. The curves are means over the 20 seeds for the breadth-first and random orderings, the backlink count $\mathrm{IB}'_U$, AWP$'$ and C-PR$'$; the shaded bands give the minimum and maximum over the seeds for $\mathrm{IB}'_U$ and C-PR$'$, the dotted line is the ideal order, and the vertical line marks $10\%$ of the universe, the range of $\mathrm{AUC}_{10\%}$. The test was defined after the registered results were known; it decides no rule of the analysis plan or the registration.}
\label{fig:discovery-curves-out}
\end{figure}

% !TEX root = ../main.tex
\section{EndorseRank and AWP Side by Side}
\label{ch:appendix-er-awp}

Chapter~\ref{ch:results} reports every score of the thesis together and reads each PageRank against the raw degree it smooths. This appendix collects the direct comparisons between the two single-layer scores, EndorseRank and AWP: their family means and paired differences in the same window, their per-proxy coefficients, and their paired differences out of window. None of these comparisons belongs to rule~A or to rule~B, and none carries a claim of the thesis. The two scores walk on different edges and are built for different constructs, so on each check of a whole cohort the score ahead is the one whose edges the check counts. Their coverage differs as well: $5{,}193$ of the matched contracts lie outside the allowance graph and $2{,}384$ outside the transfer graph, and each score is zero there (Table~\ref{tab:tie-sensitivity}). Out of window EndorseRank ranks later new approval pairs below the in-approve degree, with uniform and with activity restarts, while AWP ranks later new senders slightly above the transfer in-degree on the whole cohorts and below it among the W0 spenders that another contract had paid before the freeze (Section~\ref{sec:holdout-results}; computed after the labels were known), so a lead of EndorseRank over AWP does not show that propagation over approvals adds anything to the count of approving owners.

\dissertationSubheading[sec:er-awp-same-window]{Same window}

Table~\ref{tab:alignment-hybrid} gives the family means of the two scores with their intervals, and the last row of Table~\ref{tab:tie-sensitivity} their agreement. The two patterns cross: AWP tracks the transfer family most closely and the allowance family least, while EndorseRank tracks the allowance family most closely. Table~\ref{tab:tau-diff} tests five paired differences on the same contract resamples (Section~\ref{sec:statistical-tests}), numbered i to v in the order of its rows.

% Generated by experiments/scripts/export_latex.py -- do not edit by hand.
% Source: experiments/data/2-processed-tables-and-evaluations/revised-contract-12m/same-window/eval_summary.json
\begin{table}[htbp]
\centering
\caption[Same-window paired contrasts between EndorseRank and AWP.]{Paired bootstrap tests of differences between Kendall $\tau$ coefficients of EndorseRank and AWP in the same window (matched cohort, $n=6{,}443$ contracts). Family entries use the family-mean $\tau$; the same contract resamples are used for $a$ and $b$, so the interval is on the paired difference. An interval excluding 0 indicates a difference not attributable to sampling variation at the 5\% level. CI = 95\% percentile bootstrap over contracts (400 paired resamples, seed 42).}
\label{tab:tau-diff}
\fitwidth{%
\begin{tabular}{lccccc}
\toprule
Contrast ($a$ minus $b$) & $\tau(a)$ & $\tau(b)$ & $\Delta\tau$ & 95\% CI of $\Delta\tau$ & $P(\Delta\tau>0)$ \\
\midrule
EndorseRank allowance minus EndorseRank transfer & 0.839 & 0.297 & +0.542 & [0.520, 0.563] & 1.000 \\
AWP transfer minus EndorseRank transfer & 0.755 & 0.297 & +0.458 & [0.438, 0.476] & 1.000 \\
EndorseRank allowance minus AWP allowance & 0.839 & 0.227 & +0.613 & [0.591, 0.632] & 1.000 \\
AWP stability minus EndorseRank stability & 0.643 & 0.220 & +0.422 & [0.404, 0.438] & 1.000 \\
EndorseRank in-approve degree minus EndorseRank in-degree & 0.847 & 0.305 & +0.542 & [0.521, 0.564] & 1.000 \\
\bottomrule
\end{tabular}
}
\end{table}


All five intervals lie above zero. AWP's transfer alignment is higher than EndorseRank's (contrast~ii, $+0.458$, $[0.438,0.476]$), EndorseRank's allowance alignment is higher than AWP's (contrast~iii, $+0.613$, $[0.591,0.632]$), and AWP leads EndorseRank on the Sybil-stability family (contrast~iv, $+0.422$, $[0.404,0.438]$). The remaining two compare EndorseRank with itself: its allowance alignment is higher than its transfer alignment (contrast~i, $+0.542$, $[0.520,0.563]$), and its in-approve-degree alignment is higher than its in-degree alignment (contrast~v, $+0.542$, $[0.521,0.564]$). Ties weigh on these coefficients: EndorseRank ties $65.4\%$ of the contract pairs and AWP $13.7\%$, mostly among the contracts outside each graph (Section~\ref{sec:rank-divergence}; Appendix~\ref{sec:kendall-properties}). With those contracts kept as isolated nodes and not scored zero, EndorseRank's allowance mean rises, its transfer and Sybil-stability means fall and AWP's means stay as they are, so the gaps of contrasts~i to~iv widen in their point estimates (Table~\ref{tab:tie-sensitivity}).

The opposite directions of contrasts~ii and~iii follow from the constructs. Each score is ahead on the family built from its own edges, and the main proxies of each family include the in-degree that the corresponding PageRank smooths, so the agreement is partly built in and checks the intended construct of each score. Neither difference says which score is the better measure of reputation.

\dissertationSubheading[sec:proxy-family-detail]{Proxy detail}

Each family mean condenses several per-proxy coefficients. Tables~\ref{tab:alignment-transfer} to~\ref{tab:alignment-sybil-stability} report them one proxy at a time for EndorseRank and AWP, as Spearman $\rho$, Kendall $\tau$ and the 95\% interval of $\tau$; the transfer proxies are the ones used in the evaluation of AWP by its authors \parencite{do2023awp}. Every proxy is computed from the events between contracts, and a contract outside a layer has zero on the proxies of that layer.

% Generated by experiments/scripts/export_latex.py -- do not edit by hand.
% Source: experiments/data/2-processed-tables-and-evaluations/revised-contract-12m/same-window/eval_summary.json
\begin{table}[htbp]
\centering
\caption[Domain alignment of EndorseRank and AWP with the transfer proxies.]{Domain alignment of EndorseRank and AWP with the transfer proxies (in-degree and in-value of each contract), same window (matched cohort, $n=6{,}443$ contracts). Kendall $\tau$ and Spearman $\rho$ between scores and proxy rankings. CI = 95\% percentile bootstrap over contracts (400 paired resamples, seed 42).}
\label{tab:alignment-transfer}
\fitwidth{%
\begin{tabular}{lccc}
\toprule
Proxy & Spearman $\rho$ & Kendall $\tau$ & 95\% CI of $\tau$ \\
\midrule
\multicolumn{4}{l}{\textit{EndorseRank}} \\
In-degree & 0.345 & 0.305 & [0.287, 0.326] \\
In-value (sum inbound) & 0.338 & 0.289 & [0.269, 0.309] \\
\addlinespace
\multicolumn{4}{l}{\textit{AWP}} \\
In-degree & 0.873 & 0.764 & [0.757, 0.772] \\
In-value (sum inbound) & 0.872 & 0.746 & [0.738, 0.754] \\
\bottomrule
\end{tabular}
}
\end{table}


On the transfer family (Table~\ref{tab:alignment-transfer})\label{sub:transfer-family} AWP agrees with the in-value almost as closely as with the in-degree, for the reason given in Section~\ref{sec:same-window-alignment}. EndorseRank's transfer coefficients have intervals that overlap neither zero nor AWP's intervals. Part of them is coverage: the $2{,}173$ contracts outside both graphs lie at the bottom of both orderings, and among the $4{,}270$ contracts inside the union graph the two coefficients are $0.174$ ($[0.146,0.200]$) and $0.153$ ($[0.127,0.177]$).

% Generated by experiments/scripts/export_latex.py -- do not edit by hand.
% Source: experiments/data/2-processed-tables-and-evaluations/revised-contract-12m/same-window/eval_summary.json
\begin{table}[htbp]
\centering
\caption[Domain alignment of EndorseRank and AWP with the allowance proxies.]{Domain alignment of EndorseRank and AWP with the allowance proxies (in-approve degree and value of each spender contract), same window (matched cohort, $n=6{,}443$ contracts). Kendall $\tau$ and Spearman $\rho$ between scores and proxy rankings. CI = 95\% percentile bootstrap over contracts (400 paired resamples, seed 42).}
\label{tab:alignment-allowance}
\fitwidth{%
\begin{tabular}{lccc}
\toprule
Proxy & Spearman $\rho$ & Kendall $\tau$ & 95\% CI of $\tau$ \\
\midrule
\multicolumn{4}{l}{\textit{EndorseRank}} \\
In-approve degree & 0.883 & 0.847 & [0.835, 0.857] \\
In-approve value (sum) & 0.882 & 0.831 & [0.820, 0.841] \\
\addlinespace
\multicolumn{4}{l}{\textit{AWP}} \\
In-approve degree & 0.271 & 0.229 & [0.207, 0.248] \\
In-approve value (sum) & 0.269 & 0.225 & [0.204, 0.243] \\
\bottomrule
\end{tabular}
}
\end{table}


On the allowance family (Table~\ref{tab:alignment-allowance})\label{sub:allowance-family} EndorseRank follows the in-approve degree and the in-approve value about equally, for the reason given in Section~\ref{sec:same-window-alignment}. AWP's coefficients lie above zero, also among the contracts inside the union graph ($0.122$, $[0.091,0.151]$, and $0.116$, $[0.087,0.145]$): among the cohort's contracts, those with capped allowances from more owner contracts also tend to be central in the transfer graph.

% Generated by experiments/scripts/export_latex.py -- do not edit by hand.
% Source: experiments/data/2-processed-tables-and-evaluations/revised-contract-12m/same-window/eval_summary.json
\begin{table}[htbp]
\centering
\caption[Domain alignment of EndorseRank and AWP with the stability proxies.]{Domain alignment of EndorseRank and AWP with the stability proxies from transfer activity, same window (matched cohort, $n=6{,}443$ contracts). Kendall $\tau$ and Spearman $\rho$ between scores and proxy rankings. CI = 95\% percentile bootstrap over contracts (400 paired resamples, seed 42).}
\label{tab:alignment-sybil-stability}
\fitwidth{%
\begin{tabular}{lccc}
\toprule
Proxy & Spearman $\rho$ & Kendall $\tau$ & 95\% CI of $\tau$ \\
\midrule
\multicolumn{4}{l}{\textit{EndorseRank}} \\
Inbound counterparty ratio & 0.177 & 0.148 & [0.130, 0.166] \\
Transfer tenure (days) & 0.317 & 0.271 & [0.253, 0.289] \\
Active months & 0.281 & 0.243 & [0.224, 0.262] \\
\addlinespace
\multicolumn{4}{l}{\textit{AWP}} \\
Inbound counterparty ratio & 0.648 & 0.428 & [0.414, 0.443] \\
Transfer tenure (days) & 0.862 & 0.711 & [0.702, 0.719] \\
Active months & 0.917 & 0.789 & [0.781, 0.797] \\
\bottomrule
\end{tabular}
}
\end{table}


On the Sybil-adjusted stability family (Table~\ref{tab:alignment-sybil-stability})\label{sub:sybil-family} AWP is ahead of EndorseRank on all three proxies, as one would expect from a transfer score that rewards activity spread over the window; the paired difference on the family mean excludes zero (Table~\ref{tab:tau-diff}, contrast~iv). Coverage shapes these coefficients. The $2{,}384$ contracts outside the transfer graph have zero on all three proxies and score zero under AWP, and $2{,}173$ of them score zero under EndorseRank as well. Among the $4{,}270$ contracts inside the union graph AWP's family mean is $0.380$ ($[0.367,0.393]$) and EndorseRank's $0.028$ ($[0.009,0.046]$); on the inbound counterparty ratio AWP's coefficient does not exclude zero there ($0.018$, $[-0.007,0.042]$), and EndorseRank's is negative ($-0.091$, $[-0.112,-0.068]$). The ratio falls when the same senders return, so inside the graph the contracts that EndorseRank ranks high are, on this measure, contracts whose senders return.

\dissertationSubheading[sec:er-awp-holdout]{Out of window}

Table~\ref{tab:endorserank-awp-holdout} collects the paired differences between the two scores on the spender cohorts of both label windows. In every row $a$ is EndorseRank and $b$ AWP. The rows were computed after the labels were known.

% Generated by experiments/scripts/export_latex.py (export_latex_c12.py) -- do not edit by hand.
% Source: experiments/data/2-processed-tables-and-evaluations/revised-contract-12m/holdout-w0/eval_summary.json
% Source: experiments/data/2-processed-tables-and-evaluations/revised-contract-12m/supplementary/supplementary.json
% Source: experiments/data/2-processed-tables-and-evaluations/revised-contract-12m/window-w1/eval_summary.json
\begin{table}[htbp]
\centering
\caption[Paired bootstrap differences between EndorseRank and AWP out of window.]{Paired bootstrap differences between EndorseRank and AWP out of window (W0 spenders $n=7{,}864$; W1 spenders $n=7{,}915$). None of them enters rule A or the two applied criteria of rule B: the rows were computed after the labels were known, the W1 rows from the saved W1 scores with the resamples of Table~\ref{tab:fresh-holdout}. Within each block $a$ and $b$ share the contract resamples (400; seed 42).}
\label{tab:endorserank-awp-holdout}
\fitwidth{%
\begin{tabular}{llccccc}
\toprule
Contrast ($a$ minus $b$) & Label & $\tau(a)$ & $\tau(b)$ & $\Delta\tau$ & 95\% CI of $\Delta\tau$ & $P(\Delta\tau>0)$ \\
\midrule
\multicolumn{7}{l}{\emph{W0, spenders ($n=7{,}864$)}} \\
EndorseRank minus AWP & new approval pairs & 0.168 & 0.098 & +0.070 & [0.042, 0.096] & 1.000 \\
EndorseRank minus AWP & new transfer senders & 0.084 & 0.330 & -0.246 & [-0.271, -0.220] & 0.000 \\
\midrule
\multicolumn{7}{l}{\emph{W1, spenders ($n=7{,}915$)}} \\
EndorseRank minus AWP & new approval pairs & 0.147 & 0.064 & +0.083 & [0.059, 0.110] & 1.000 \\
EndorseRank minus AWP & new transfer senders & 0.064 & 0.314 & -0.250 & [-0.273, -0.225] & 0.000 \\
\bottomrule
\end{tabular}
}
\end{table}


On the spenders of W0, EndorseRank ranks new approval pairs above AWP, a paired difference of $+0.070$ (interval $[0.042,0.096]$). On new transfer senders the order reverses, with a difference of $-0.246$ ($[-0.271,-0.220]$). Every label falls after all scored events. The first gap rests on coverage. AWP scores zero for the $2{,}421$ spenders without a transfer pair before the freeze, nearly a third of the cohort, which only allowance edges can order, and EndorseRank for the $5{,}966$ without an allowance edge of positive weight. On the $4{,}569$ spenders that another contract had paid before the freeze, AWP's coefficient is the higher on both labels (Table~\ref{tab:holdout-receiving}), without a paired test.

In July--September the pattern repeats: EndorseRank leads on new approval pairs ($+0.083$, $[0.059,0.110]$) and AWP on new senders ($-0.250$, $[-0.273,-0.225]$).

% !TEX root = ../main.tex
\section{The Registered Analysis}
\label{sec:map-registered}

The registered analysis is the one that the analysis plan and the registration of W1 fixed before the study extracted the data it judges (Appendix~\ref{sec:pipeline-stages}): amounts in raw token base units, every ERC-20 token, the value transform with $b=1$, scores from every event since 1~October~2023, graphs that hold every account of the cohort's events, and the plan's cohort of $27{,}844$ contracts. After the registered results were known, the revised analysis was defined: amounts in US dollars on five tokens; scores from the twelve months before each freeze; graphs whose nodes are contracts other than smart accounts, with edges read from the whole chain; each allowance counted up to its owner's balance at the freeze, with balances reconstructed from \texttt{Transfer} logs; constants set on the contract graph; and the discovery-order test. The chapters report this analysis (Section~\ref{sec:amount-revision}), in which rules~A and~B are applied to the scores of the contract graph as a re-analysis. Table~\ref{tab:registered-vs-revised} sets the decision contrasts and verdicts of the registered analysis beside the same contrasts on those scores; its last four rows, degree contrasts computed after the labels were known, carry no verdict. Of rule~B the thesis applies two of its registered criteria, F1 and F2; under the registration the rule holds only if each of its criteria holds. The registration is \texttt{docs/registration\_w1.md} at commit \texttt{1747e64}. The two columns reach the same verdicts: neither part of rule~A is met, F1 holds and F2 fails, also in the two registered sensitivity analyses, so rule~B is not met, and F6 fails. Only the verdicts of the registered column are outcomes of tests fixed before the study extracted the data they judge. The complete registered results, with every coefficient, interval and verdict, are published unchanged in \texttt{data/3-published-results-for-thesis/} with their digest, \texttt{results\_digest.md}.

% Generated by experiments/scripts/export_latex.py (export_latex_c12.py) -- do not edit by hand.
% Source: experiments/data/2-processed-tables-and-evaluations/same-window/eval_summary.json
% Source: experiments/data/2-processed-tables-and-evaluations/holdout-w0/eval_summary.json
% Source: experiments/data/2-processed-tables-and-evaluations/registered-w1/eval_summary.json
% Source: experiments/data/2-processed-tables-and-evaluations/revised-contract-12m/same-window/eval_summary.json
% Source: experiments/data/2-processed-tables-and-evaluations/revised-contract-12m/holdout-w0/eval_summary.json
% Source: experiments/data/2-processed-tables-and-evaluations/revised-contract-12m/window-w1/eval_summary.json
\begin{table}[htbp]
\centering
\caption[Decision contrasts of the registered analysis and of the contract-graph analysis.]{Decision contrasts and verdicts of the registered analysis (amounts in raw base units, every ERC-20 token, $b=1$, scores from every event since 1~October~2023, graphs that hold every account of the cohort's events), exactly as obtained under the analysis plan and the W1 registration, beside the same contrasts on the scores of the contract graph (amounts in USD, the five tokens, scores from the twelve months before each freeze, nodes that are contracts other than smart accounts, each allowance counted up to its owner's balance). The contract-graph analysis was defined after the registered results were known; its column applies the registered rules and thresholds to its scores and is not a registered test. Under the registration rule B holds only if each of its criteria holds; the table applies two of them, F1 and F2, in both columns. The last block holds degree contrasts that no rule of either analysis decides; they are reported in both. $\Delta\tau$ with 95\% paired bootstrap interval (400 resamples, seed 42); the rules and their thresholds are in Sections~\ref{sec:holdout-protocol} and~\ref{sub:fresh-holdout}.}
\label{tab:registered-vs-revised}
\fitwidth{%
\begin{tabular}{lcc}
\toprule
Quantity & Registered analysis & Contract-graph analysis \\
\midrule
Matched cohort (contracts) & 27{,}844 & 6{,}443 \\
W0 spender cohort & 33{,}101 & 7{,}864 \\
W1 spender cohort & 34{,}587 & 7{,}915 \\
EndorseRank vs.\ AWP, same window ($\tau$) & 0.217 [0.209, 0.225] & 0.275 [0.258, 0.294] \\
\midrule
A, W0: C-PR minus C-PR ($\lambda=1$), new approval pairs & $+0.008$ [0.003, 0.013] & $-0.025$ [-0.045, -0.004] \\
A, W0: C-PR minus C-PR ($\lambda=0$), new transfer senders & $-0.030$ [-0.034, -0.026] & $-0.014$ [-0.019, -0.008] \\
Rule A, primary part & not met & not met \\
Rule A, secondary part & not met & not met \\
\midrule
B, F1: C-PR minus C-PR ($\lambda=0$), new approval pairs & $+0.101$ [0.092, 0.111] & $+0.074$ [0.060, 0.087] \\
B, F2: C-PR minus C-PR ($\lambda=0$), new transfer senders & $-0.025$ [-0.029, -0.021] & $-0.016$ [-0.022, -0.012] \\
F1 and F2, two of rule B's registered criteria & not met (F2 fails) & not met (F2 fails) \\
F1 and F2, sensitivity 1 (AWP as comparator) & not met (F2 fails) & not met (F2 fails) \\
F1 and F2, sensitivity 2 (isolated nodes) & not met (F2 fails) & not met (F2 fails) \\
F6 (rule A, primary part, on W1) & fails & fails \\
\midrule
W0: EndorseRank minus in-approve degree, new pairs & $-0.011$ [-0.016, -0.006] & $-0.036$ [-0.046, -0.027] \\
W0: AWP minus transfer in-degree, new senders & $+0.038$ [0.035, 0.042] & $+0.012$ [0.003, 0.023] \\
W1: EndorseRank minus in-approve degree, new pairs & $-0.029$ [-0.032, -0.025] & $-0.034$ [-0.043, -0.025] \\
W1: AWP minus transfer in-degree, new senders & $+0.020$ [0.016, 0.023] & $+0.011$ [0.001, 0.022] \\
\bottomrule
\end{tabular}
}
\end{table}


% !TEX root = ../main.tex
\section{Ethics Clearance Certificate}
\label{ch:appendix-ethics}

\dissertationSubheading[sec:ethics-certificate]{Certificate}

This study uses public blockchain data only and involves no human participants (Section~\ref{sec:ethics}). The ethics clearance certificate issued by the Unisa School of Computing / College of Agriculture and Environmental Sciences research ethics review committee is reproduced below.

% TODO (user action): replace the placeholder box with the certificate PDF once
% the reference number, date and scanned document are available, e.g.
%   \includegraphics[width=\textwidth,page=1]{Figures/ethics-clearance.pdf}
% Do not invent a reference number or date.
\begin{center}
\fbox{\parbox{0.9\textwidth}{\centering\vspace{2em}
Ethics clearance certificate\\[0.5em]
Reference: [Unisa SoC/CAES ethics clearance ref: TBD]\\
Date of issue: [TBD]\\[0.5em]
(Scanned certificate to be inserted here.)
\vspace{2em}}}
\end{center}

